\subsection*{哪些推理其实没有用到箭头形状}
回顾前面，许多证明只用了相加和数乘，并没有使用长度或角度。若另一类对象也支持相同规则，就可以搬用这些推理。函数相加按每个输入逐点相加，数乘把每个输出同比例缩放；多项式与矩阵也一样。

所有次数不超过二的实多项式，都唯一写成 $c_0+c_1x+c_2x^2$。所以 $1,x,x^2$ 扮演三根基，系数是坐标。这里“空间”首先是一种组织对象和操作的结构，不必真的把全部函数画进三维图中。

\subsection*{求导也能写成矩阵}
求导把三根基依次送到 $0,1,2x$。把结果在原基下的坐标排成列，得到下面的求导矩阵。计算多项式导数就变成系数向量的矩阵变换。常数在核中，因为求导看不见它们；输出最多一次，所以像空间也受到限制。

这说明抽象并非远离应用：箭头上的基、核、线性映射，直接解释了函数运算。傅里叶级数会换用正弦和余弦作基本方向，热方程会研究各方向随时间的倍率，后续章节由此连接。

迁移结构时仍要核对封闭性。“恰好二次”的多项式相减可能成为一次或零，所以不组成这种向量空间；任意函数也不都可导，不能把求导当成整个函数空间上的算子。抽象的工作，是保留推理需要的条件，并明确哪些对象真正满足它们。
